\documentclass[a4paper, serif, 10pt]{moderncv}

%% moderncv
% color and style
\moderncvcolor{black}
\moderncvstyle{banking}

% renew moderncv command to adapt for CJK colon
\renewcommand*{\cvitem}[3][.25em]{%
  \ifstrempty{#2}{}{\hintstyle{#2}：}{#3}%
  \par\addvspace{#1}}

\renewcommand*{\cvitemwithcomment}[4][.25em]{%
  \savebox{\cvitemwithcommentmainbox}{\ifstrempty{#2}{}{\hintstyle{#2}：}#3}%
  \setlength{\cvitemwithcommentmainlength}{\widthof{\usebox{\cvitemwithcommentmainbox}}}%
  \setlength{\cvitemwithcommentcommentlength}{\maincolumnwidth-\separatorcolumnwidth-\cvitemwithcommentmainlength}%
  \begin{minipage}[t]{\cvitemwithcommentmainlength}\usebox{\cvitemwithcommentmainbox}\end{minipage}%
  \hfill% fill of \separatorcolumnwidth
  \begin{minipage}[t]{\cvitemwithcommentcommentlength}\raggedleft\small\itshape#4\end{minipage}%
  \par\addvspace{#1}}

%% page layout/margins
\usepackage[top=2.5cm, bottom=2.5cm, left=1.5cm, right=1.5cm]{geometry}

\nopagenumbers{}

%% Babel config for English language
\usepackage[english]{babel}

%% fontspec
\usepackage{fontspec}

\IfFontExistsTF{Linux Libertine}{
  \setmainfont[Ligatures={TeX, Common}, Numbers=Lining, ItalicFont=Linux Libertine]{Linux Libertine}
}{}
\IfFontExistsTF{Linux Libertine O}{
  \setmainfont[Ligatures={TeX, Common}, Numbers=Lining, ItalicFont=Linux Libertine O]{Linux Libertine O}
}{}

%% CTeX
% CJK support, used to show CJK characters in the resume
%
% - heading=false: disable ctex heading
% - punct=kaiming: use kaiming punctuations styles for CJK
% - scheme=plain: use plain scheme, do not override \normalsize font size
% - space=auto: space settings for CJK characters
%
% CJK fonts are resolved per locale language, e.g. Noto Serif/Sans CJK JP
% for Japanese resumes, with fallback to the next region if not installed
%
% ref:
% - http://ctan.mirrorcatalogs.com/language/chinese/ctex/ctex.pdf
\usepackage[UTF8, heading=false, punct=kaiming, scheme=plain, space=auto]{ctex}

\IfFontExistsTF{Noto Serif CJK TC}{
  \setCJKmainfont{Noto Serif CJK TC}
}{}
\IfFontExistsTF{Noto Sans CJK TC}{
  \setCJKsansfont{Noto Sans CJK TC}
}{}

%% line spacing
\usepackage{setspace}
\setstretch{1.125}

%% URL styling - use normal text instead of monospace
\urlstyle{same}

%% Patch \httplink and \httpslink to handle full URLs
% The original moderncv macros always prepend http:// or https:// to URLs.
% This causes issues when the URL already has a protocol (e.g., https://example.com)
% resulting in malformed URLs like https://https://example.com.
% This patch checks if the URL already contains "://" and uses it directly if so.
% Note: We use \str_set:Nx (x-expansion) to expand macros like \@homepage first.
\ExplSyntaxOn
\str_new:N \g_yamlresume_url_str

\RenewDocumentCommand{\httpslink}{O{}m}{
  \str_set:Nx \g_yamlresume_url_str {#2}
  \str_if_in:NnTF \g_yamlresume_url_str {://}
    { \url{#2} }
    { \url{https://#2} }
}

\RenewDocumentCommand{\httplink}{O{}m}{
  \str_set:Nx \g_yamlresume_url_str {#2}
  \str_if_in:NnTF \g_yamlresume_url_str {://}
    { \url{#2} }
    { \url{http://#2} }
}
\ExplSyntaxOff

\name{林雅婷}{}
\title{註冊護理師｜急重症與急診照護}
\phone[mobile]{+886 912 345 678}
\email{yating.lin.rn@example.com}
\homepage{https://www.linkedin.com/in/yating-lin-rn}

\address{中國臺灣，台北市，台北市大安區仁愛路四段123號，106}{}{}

\extrainfo{{\small \faLinkedin}\ \href{https://www.linkedin.com/in/yating-lin-rn}{@yating-lin-rn} {} {} {} • {} {} {} 
{\small \faFacebook}\ \href{https://www.facebook.com/yating.lin.rn}{@yating.lin.rn}}

\begin{document}

\maketitle

\section{簡介}

\cvline{}{具備六年以上急診與重症加護臨床經驗的註冊護理師，熟悉檢傷分類、生命徵象監測、靜脈注射、傷口照護及病患衛教。擅長在高壓環境中進行快速評估與跨團隊溝通，致力於提供以病人為中心的高品質照護。
%
曾參與急診流程優化與譫妄預防專案，具備臨床教學與品質改善經驗。持續進修進階護理與實證照護知能，期望於急重症單位發揮專業，提升病人安全與照護品質。}
\section{教育背景}

\cventry{2014年9月–2018年6月}
        {學士，護理學，成績：3.9/4.0}
        {國立臺灣大學醫學院護理學系}
        {\url{https://www.mc.ntu.edu.tw/nursing}}
        {}
        {\begin{itemize}
\item 以優異成績完成護理學學士學位，並通過全國護理師高考
\item 於臺大醫院完成臨床實習，累積內外科、急診與加護病房照護經驗
\item 擔任護理系學會幹部，協助推動校內外健康促進活動
\end{itemize}
\textbf{課程}：基本護理學、內外科護理學、急重症護理學、產科護理學、兒科護理學、精神科護理學、社區護理學、護理倫理與法律}
\section{工作經歷}

\cventry{2018年7月–2022年6月}
        {註冊護理師}
        {臺大醫院急診部}
        {\url{https://www.ntuh.gov.tw}}
        {}
        {\begin{itemize}
\item 負責急診病患檢傷分類、生命徵象監測與初步處置，每日平均照護 20 至 30 位病患
\item 執行靜脈注射、抽血、傷口換藥、導尿及心電圖監測等護理技術
\item 與急診醫師、專科護理師及社工師協作，執行跨團隊照護與轉院/住院安排
\item 提供病患及家屬疾病衛教、出院準備與社區資源轉介
\item 參與急診壅塞改善專案，協助優化檢傷流程與縮短候診時間
\end{itemize}
\textbf{關鍵字}：急診護理、檢傷分類、靜脈注射、病患衛教、跨團隊照護}

\cventry{2022年7月至今}
        {資深註冊護理師}
        {臺北榮民總醫院重症加護病房}
        {\url{https://www.vghtpe.gov.tw}}
        {}
        {\begin{itemize}
\item 照護呼吸衰竭、敗血症、急性腎損傷及術後重症病患，熟悉呼吸器、血液透析與葉克膜輔助照護
\item 執行重症病患之持續性生命徵象監測、給藥、輸液與侵入性管路照護
\item 帶領新進護理師與實習生，規劃臨床教學與技術考核
\item 參與加護病房譫妄評估與預防流程改善，降低譫妄發生率
\item 與醫師、藥師、營養師及復健師合作，執行每日跨團隊查房與照護計畫
\end{itemize}
\textbf{關鍵字}：重症照護、呼吸器、血液透析、臨床教學、品質改善}
\section{語言}

\cvline{中文}{母語或雙語水平 \hfill \textbf{關鍵字}：母語、國語}
\cvline{英語}{完全專業水平 \hfill \textbf{關鍵字}：TOEIC 850、護理英語}
\section{專業技能}

\cvline{臨床護理}{專家 \hfill \textbf{關鍵字}：檢傷分類、生命徵象監測、靜脈注射、傷口照護、管路照護}
\cvline{急重症照護}{高級 \hfill \textbf{關鍵字}：呼吸器照護、血液透析、葉克膜輔助、敗血症照護、急救處置}
\cvline{病患溝通與衛教}{高級 \hfill \textbf{關鍵字}：醫病溝通、出院準備、家屬支持、健康促進}
\cvline{醫療資訊系統}{中級 \hfill \textbf{關鍵字}：HIS、電子病歷、醫囑系統}
\cvline{品質管理與教學}{高級 \hfill \textbf{關鍵字}：品質改善、臨床教學、實證護理、專案管理}
\section{榮譽}

\cventry{2021年12月}
        {臺大醫院}
        {優良護理人員獎}
        {}
        {}
        {表揚在急診照護、病人安全及團隊合作上表現優異之護理人員。}

\cventry{2023年11月}
        {臺北榮民總醫院}
        {臨床教學優良獎}
        {}
        {}
        {肯定於新進護理師臨床教學與技能指導之卓越貢獻。}
\section{證書}

\cventry{2018年9月}
        {考選部}
        {高考護理師證書}
        {\url{https://www.moex.gov.tw}}
        {}
        {}

\cventry{2020年12月}
        {台灣護理學會}
        {進階護理師 N2 證書}
        {\url{https://www.twna.org.tw}}
        {}
        {}

\cventry{2022年3月}
        {台灣急救加護醫學會}
        {高級心臟救命術 ACLS}
        {\url{https://www.seccm.org.tw}}
        {}
        {}

\cventry{2024年1月}
        {台灣急救加護醫學會}
        {基本生命支持 BLS}
        {\url{https://www.seccm.org.tw}}
        {}
        {}
\section{出版刊物}

\cventry{2022年6月}
        {運用跨團隊照護改善急診檢傷分類流程之專案}
        {台灣急診醫學會雜誌}
        {\url{https://www.sem.org.tw}}
        {}
        {\begin{itemize}
\item 探討急診檢傷流程瓶頸並提出改善策略
\item 透過跨團隊合作與資訊系統優化，縮短病患候診時間
\item 提升檢傷分類正確率與病人滿意度
\end{itemize}}
\section{項目}

\cventry{2023年1月–2023年12月}
        {針對加護病房重症病患推動譫妄評估與預防專案，透過標準化流程、教育訓練及跨團隊合作，提升護理人員對譫妄的識別與處置能力，進而改善病患預後。}
        {加護病房譫妄評估與預防流程改善}
        {\url{https://www.vghtpe.gov.tw}}
        {}
        {\begin{itemize}
\item 導入 CAM-ICU 評估工具，建立標準化譫妄篩檢流程
\item 規劃護理人員教育訓練與稽核機制
\item 結合早期活動與睡眠促進措施，降低譫妄發生率
\end{itemize}
\textbf{關鍵字}：譫妄、加護病房、品質改善、病人安全}
\section{興趣愛好}

\cvline{專業進修}{實證護理、急重症醫學、護理教育}
\cvline{休閒活動}{馬拉松、閱讀、登山}
\section{志願服務}

\cventry{2019年3月至今}
        {急救教育志工}
        {台灣紅十字會}
        {\url{https://www.redcross.org.tw}}
        {}
        {\begin{itemize}
\item 協助推廣社區急救教育，教授 CPR 與 AED 使用
\item 參與偏鄉義診與健康篩檢活動
\item 協助大型活動醫護站輪值與緊急處置
\end{itemize}}
    
\end{document}